%%%%%%%%%%%%%%%%%%%%%%%%%%%%%%%%%%%%%%%%%%%%%%%%%%%%%%%%%%%%%%%%%%%%%%%%%%%%%%%
%  Beispielfolien -- Schaufenster aller Folienbausteine
%
%  Dieses File dient als Nachschlagewerk beim Bauen der eigenen Folien und
%  gehoert NICHT in die fertige Praesentation. Vor dem Vortrag einfach den
%  \input in presentation.tex loeschen.
%
%  Uebersicht der Bausteine:
%    \EITHeadlineFT{Titel}{Inhalt}     Standardfolie mit Ueberschrift
%    \placetextbox{x}{y}{breite}{...}  frei positionierter Textblock
%    \placepictureandtext{...}         Bild mit Beschriftung darunter
%    \picdims[opt]{breite}{hoehe}{..}  Bild auf feste Masse zuschneiden
%%%%%%%%%%%%%%%%%%%%%%%%%%%%%%%%%%%%%%%%%%%%%%%%%%%%%%%%%%%%%%%%%%%%%%%%%%%%%%%

%--- Blindtext nur fuer diese Beispielfolien ----------------------------------
\newcommand{\germanlipsum}{%
	Lorem ipsum ist ein pseudo-lateinischer Text für Webdesign, Typografie,
	Layout und Printmedien. Er ersetzt Deutsch, um Designelemente gegenüber
	dem Inhalt hervorzuheben, hat also die Funktion eines Platzhalters im
	Layout, um dem Betrachter eine vorläufige Vorstellung der endgültigen
	Fassung zu vermitteln. Buchstaben, Worte, Wort- und Satzlängen geben
	ungefähr das Schriftbild eines lateinischen Textes wieder. Um die
	Aufmerksamkeit ausschließlich auf das Schriftbild zu lenken, wird ein
	zufälliger Sinn durch Wort- oder Satzkombinationen vermieden.%
}
\newcommand{\germanlipsumshort}{%
	Lorem ipsum ist ein pseudo-lateinischer Text für Webdesign, Typografie,
	Layout und Printmedien.%
}


%==============================================================================
\section{Textfolien}
%==============================================================================

\EITHeadlineFT{Folie mit Fließtext}{
	\germanlipsum
}

% \foottext wirkt bis zum Ende der umgebenden Gruppe -- die geschweiften
% Klammern um diesen Block begrenzen die Wirkung also auf diese eine Folie.
% Fuer die gesamte Praesentation stattdessen \foottext in presentation.tex
% einkommentieren.
{
\newcommand{\foottext}{Hochschule Karlsruhe --- Fakultät für Elektro- und Informationstechnik --- Vorname Nachname}
\EITHeadlineFT{Folie mit eigener Fußzeile}{
	Ist \textbf{\textbackslash foottext} definiert, ersetzt eine zentrierte
	Zeile die Standardfußzeile -- so wie auf dieser Folie.

	Ohne diese Definition stehen links Hochschule, Datum und Foliennummer.
}
}

\EITHeadlineFT{Folie mit Liste}{
	\begin{itemize}
		\item \germanlipsumshort
		\item \germanlipsumshort
		\item \germanlipsumshort
		\item \germanlipsumshort
	\end{itemize}
}


%==============================================================================
\section{Formeln}
%==============================================================================

\EITHeadlineFT{Folie mit Formel}{
	Verwiesen wird mit \textbf{\textbackslash cref}, das Wort davor kommt
	automatisch: \cref{eq:beispiel}.

	\begin{equation}
		f(x)=\left\{\begin{aligned}
			& \mathrm{e}^{|x|} & & \text{falls $|x-x^0|\leq 1/2$} \\
			& 0                & & \text{falls $|x-x^0| > 1/2$}
		\end{aligned}\right.
		\label{eq:beispiel}
	\end{equation}
}


%==============================================================================
\section{Mehrspaltige Folien}
%==============================================================================

\EITHeadlineFT{Folie mit zwei Spalten}{
	\placetextbox{\leftmarginpos}{\topmargincontent}{\columnmaxwidth}{
		\germanlipsum
	}
	\placetextbox{\secondcolumnstart}{\topmargincontent}{\columnmaxwidth}{
		\germanlipsum
	}
}

\EITHeadlineFT{Frei positionierter Text}{
	% Die Masse lassen sich als Ausdruck schreiben, TikZ rechnet das aus.
	\placetextbox{\leftmarginpos + 3cm}{\topmargincontent - 1cm}{\columnmaxwidth}{
		\germanlipsum
	}
}


%==============================================================================
\section{Tabellen}
%==============================================================================

\EITHeadlineFT{Folie mit Tabelle}{

	Verwiesen wird auf \cref{tab:messwerte}.

	\begin{table}[!h]
		\centering
		\begin{threeparttable}
			% Spaltentyp S richtet an der Kommastelle aus (siunitx).
			\begin{tabular}{@{} l S[table-format=6.0] l S[table-format=5.0] cc @{}}
				\toprule
				{Nr.} & {Kennnummer} & {Bezeichnung} & {Messwert} & \multicolumn{2}{c@{}}{Abweichung} \\
				\cmidrule(l){5-6}
				      &              &               & {(mV)}     & {unten} & {oben} \\
				\midrule
				1 & 234523 & Probe A & 234   & 0 & 0 \\
				2 & 2345   & Probe B & 2365  & 0 & 0 \\
				3 & 3453   & Probe C & 45634 & 0 & 0 \\
				4 & 83452  & Probe D & 2456  & 0 & 0 \\
				\addlinespace
				5 & 210    & Probe E & 245   & 0 & 0 \\
				6 & 3417   & Probe F & 45634 & 0 & 0 \\
				7 & 4345   & Probe G & 3456  & 0 & 0 \\
				8 & 4334   & Probe H & 3456  & 0 & 0 \\
				\bottomrule
			\end{tabular}
			\caption{Messwerte der acht Proben aus der Vergleichsmessung}
			\label{tab:messwerte}
		\end{threeparttable}
	\end{table}
}

\EITHeadlineFT{Frei positionierte Tabelle}{
	\placetextbox{0.2\paperwidth}{-0.2\paperheight}{\maximumwidth}{
		\begin{table}[!h]
			\centering
			\begin{threeparttable}
				\begin{tabular}{@{} l S[table-format=6.0] l S[table-format=5.0] cc @{}}
					\toprule
					{Nr.} & {Kennnummer} & {Bezeichnung} & {Messwert} & \multicolumn{2}{c@{}}{Abweichung} \\
					\cmidrule(l){5-6}
					      &              &               & {(mV)}     & {unten} & {oben} \\
					\midrule
					1 & 234523 & Probe A & 234   & 0 & 0 \\
					2 & 2345   & Probe B & 2365  & 0 & 0 \\
					3 & 3453   & Probe C & 45634 & 0 & 0 \\
					4 & 83452  & Probe D & 2456  & 0 & 0 \\
					\bottomrule
				\end{tabular}
				\caption{Auszug der Messwerte, frei auf der Folie platziert}
				\label{tab:messwerte-frei}
			\end{threeparttable}
		\end{table}
	}
}


%==============================================================================
\section{Abbildungen}
%==============================================================================

\EITHeadlineFT{Abbildung mit Beschriftung}{
	Verwiesen wird auf \cref{fig:zugeschnitten}.

	\begin{figure}[h]
		\centering
		\begin{measuredfigure}
			\picdims[width=\paperwidth]{\columnmaxwidth}{\maximumheight}{chapter_image_1.jpg}
			\caption{Auf feste Maße zugeschnittenes Bild}
		\end{measuredfigure}
		\MeasuredFigureLabel{fig:zugeschnitten}
	\end{figure}
}

\EITHeadlineFT{Abbildung in fester Breite}{
	Verwiesen wird auf \cref{fig:festebreite}.

	\begin{figure}[h]
		\centering
		\begin{measuredfigure}
			% Bewusst width statt scale: scale haengt an der Pixelaufloesung
			% der Bilddatei, ein Bild mit doppelter Aufloesung wuerde bei
			% gleichem scale-Wert ueber den Folienrand laufen.
			\includegraphics[width=0.45\paperwidth]{chapter_image_1.jpg}
			\caption{Bild mit fest vorgegebener Breite}
		\end{measuredfigure}
		\MeasuredFigureLabel{fig:festebreite}
	\end{figure}
}


%==============================================================================
\section{Bildfolien}
%==============================================================================

\EITHeadlineFT{Folie mit einem Bild}{
	\placetextbox{\leftmarginpos}{\topmargincontent}{\maximumwidth}{
		\picdims[width=\maximumwidth]{\maximumwidth}{\maximumheight}{chapter_image_1.jpg}
	}
}

\EITHeadlineFT{Folie mit zwei Bildern}{
	\placetextbox{\leftmarginpos}{\topmargincontent}{\columnmaxwidth}{
		\picdims[width=\paperwidth]{\columnmaxwidth}{\maximumheight}{chapter_image_1.jpg}
	}
	\placetextbox{\secondcolumnstart}{\topmargincontent}{\columnmaxwidth}{
		\picdims[width=\paperwidth]{\columnmaxwidth}{\maximumheight}{chapter_image_2.jpg}
	}
}

\EITHeadlineFT{Frei positioniertes Bild}{
	\placetextbox{\leftmarginpos + 0.2\paperwidth}{\topmargincontent - 5cm}{\maximumwidth}{
		\includegraphics[width=0.45\paperwidth]{chapter_image_1.jpg}
	}
}

\EITHeadlineFT{Frei positioniertes Bild mit Beschriftung}{
	\placepictureandtext{\leftmarginpos + 0.2\paperwidth}{\topmargincontent - 5cm}{\maximumwidth}{0cm}{0.1cm}
	{
		\includegraphics[width=0.45\paperwidth]{chapter_image_1.jpg}
	}
	{
		\color{LegendColor}\fontsize{7.7pt}{3pt}\selectfont
		Bildbeschreibung
	}
}

\EITHeadlineFT{Text links, Bild rechts}{
	\placetextbox{\leftmarginpos}{\topmargincontent}{\columnmaxwidth}{
		\germanlipsum
	}
	\placetextbox{\secondcolumnstart}{\topmargincontent}{\columnmaxwidth}{
		\picdims[width=\paperwidth]{\columnmaxwidth}{\maximumheight}{chapter_image_1.jpg}
	}
}

% Randlose Bildfolie: leere Ueberschrift, Bild ueber die volle Folienflaeche.
\EITHeadlineFT{}{
	\placetextbox{0cm}{0cm}{\paperwidth}{
		\picdims[width=\paperwidth]{\paperwidth}{16cm}{chapter_image_1.jpg}
	}
}


%==============================================================================
\section{Quellen}
%==============================================================================

\EITHeadlineFT{Literatur}{
	\begin{itemize}
		\item \fullcite{article_key}
		\item \fullcite{book_key}
	\end{itemize}
}
